\begin{examplebox}
	\textbf{\textcolor{CExample}{例 1（基础）}}

	\textbf{\textcolor{CExample}{题目：}}
	计算 $\sin75^\circ + \sin15^\circ$。

	\textbf{\textcolor{CExample}{解题步骤：}}
	\begin{description}[style=nextline, leftmargin=2.8em, labelsep=0.5em, itemsep=0.45em]
		\item[\textbf{第一步（识别结构）：}] 原式为两正弦之和，符合公式 $\sin\alpha+\sin\beta$ 形式。
			\par\textbf{理由：} 可套用和差化积公式简化。
		\item[\textbf{第二步（代入公式）：}] 取 $\alpha=75^\circ,\ \beta=15^\circ$，则
			\[
				\sin75^\circ+\sin15^\circ = 2\sin\frac{75^\circ+15^\circ}{2}\cos\frac{75^\circ-15^\circ}{2}.
			\]
			\par\textbf{理由：} 直接运用 $\sin\alpha+\sin\beta = 2\sin\frac{\alpha+\beta}{2}\cos\frac{\alpha-\beta}{2}$。
		\item[\textbf{第三步（计算数值）：}]
			\[
				\begin{aligned}
					2\sin45^\circ\cos30^\circ &= 2\cdot\frac{\sqrt{2}}{2}\cdot\frac{\sqrt{3}}{2} \\
					&= \frac{\sqrt{6}}{2}.
			\end{aligned}
			\]
			\par\textbf{理由：} 代入特殊角的三角函数值得结果。
	\end{description}

	\textbf{\textcolor{CExample}{关键结论：}}
	$\sin75^\circ+\sin15^\circ = \dfrac{\sqrt{6}}{2}$。

	\medskip
	\textbf{\textcolor{CExample}{例 2（稍变形）}}

	\textbf{\textcolor{CExample}{题目：}}
	计算 $\cos105^\circ - \cos75^\circ$。

	\textbf{\textcolor{CExample}{解题步骤：}}
	\begin{description}[style=nextline, leftmargin=2.8em, labelsep=0.5em, itemsep=0.45em]
		\item[\textbf{第一步（识别结构）：}] 原式为两余弦之差，符合公式 $\cos\alpha-\cos\beta$ 形式。
			\par\textbf{理由：} 可套用和差化积公式。
		\item[\textbf{第二步（代入公式）：}] 取 $\alpha=105^\circ,\ \beta=75^\circ$，则
			\[
				\begin{aligned}
					\cos105^\circ-\cos75^\circ &= -2\sin\frac{105^\circ+75^\circ}{2}\sin\frac{105^\circ-75^\circ}{2} \\
					&= -2\sin90^\circ\sin15^\circ.
			\end{aligned}
			\]
			\par\textbf{理由：} 使用 $\cos\alpha-\cos\beta = -2\sin\frac{\alpha+\beta}{2}\sin\frac{\alpha-\beta}{2}$。
		\item[\textbf{第三步（化简求值）：}]
			\[
				\begin{aligned}
					-2\cdot1\cdot\sin15^\circ &= -2\sin15^\circ \\
					&= -2\cdot\frac{\sqrt{6}-\sqrt{2}}{4} \\
					&= -\frac{\sqrt{6}-\sqrt{2}}{2} \\
					&= \frac{\sqrt{2}-\sqrt{6}}{2}.
			\end{aligned}
			\]
			\par\textbf{理由：} $\sin90^\circ=1$，$\sin15^\circ=\frac{\sqrt{6}-\sqrt{2}}{4}$，代入得结果。
	\end{description}

	\textbf{\textcolor{CExample}{关键结论：}}
	$\cos105^\circ-\cos75^\circ = \dfrac{\sqrt{2}-\sqrt{6}}{2}$。

	\medskip
	\textbf{\textcolor{CExample}{例 3（综合应用）}}

	\textbf{\textcolor{CExample}{题目：}}
	证明 $\sin^2 A - \sin^2 B = \sin(A+B)\sin(A-B)$。

	\textbf{\textcolor{CExample}{解题步骤：}}
	\begin{description}[style=nextline, leftmargin=2.8em, labelsep=0.5em, itemsep=0.45em]
		\item[\textbf{第一步（平方差分解）：}] 左边 $\sin^2 A-\sin^2 B = (\sin A-\sin B)(\sin A+\sin B)$。
			\par\textbf{理由：} 平方差公式。
		\item[\textbf{第二步（和差化积）：}] 对 $\sin A\pm\sin B$ 分别化积：
			\[
				\sin A-\sin B = 2\cos\frac{A+B}{2}\sin\frac{A-B}{2},\quad
				\sin A+\sin B = 2\sin\frac{A+B}{2}\cos\frac{A-B}{2}.
			\]
			\par\textbf{理由：} 应用正弦和差化积公式。
		\item[\textbf{第三步（合并化简）：}]
			\[
				\begin{aligned}
					(\sin A-\sin B)(\sin A+\sin B)
					&= \left(2\cos\frac{A+B}{2}\sin\frac{A-B}{2}\right)\left(2\sin\frac{A+B}{2}\cos\frac{A-B}{2}\right) \\
					&= 4\sin\frac{A+B}{2}\cos\frac{A+B}{2}\sin\frac{A-B}{2}\cos\frac{A-B}{2} \\
					&= \left(2\sin\frac{A+B}{2}\cos\frac{A+B}{2}\right)\left(2\sin\frac{A-B}{2}\cos\frac{A-B}{2}\right) \\
					&= \sin(A+B)\sin(A-B).
			\end{aligned}
			\]
			\par\textbf{理由：} 利用二倍角公式 $2\sin\theta\cos\theta = \sin2\theta$ 得到最终结果。
	\end{description}

	\textbf{\textcolor{CExample}{关键结论：}}
	恒等式成立。
\end{examplebox}
